%%%%%%%%%%%%%%%%%%%%%%%%%%%%%%%%%%%%%%%%%%%%%%%%%%%%%%%%%%%%%%%%%%%%%%%%%%%%%%%%
%%
\cleardoubleoddpage%  Make sure to start each chapter on a new odd page
\chapter{Conclusion}
\label{ch:conclusion}

% Drafted 2026-09-30 by an external LLM from prompts/PROMPT_final_chapters.md
% and audited line by line against that fact sheet (see NOTES.md).

This thesis set out to answer one question: \emph{What is needed to
forecast the progression of Geographic Atrophy from longitudinal OCT?} It
was split into two parts: (1) What kind of framework does the forecast
need? (2) What kind of operator does that framework need?

For the first part, the forecast requires a setup that accommodates the
irregular visit schedules of clinical data and treats every candidate
model equally. The framework of Chapter~\ref{ch:method} implements these
requirements with a residual update conditioned on the elapsed time
$\Delta t$, a zero-initialised output so that every model starts from
persistence, a training curriculum measured in elapsed days, and one loss
and one evaluation for all models. The same framework served operators
from unrelated literatures without modification, and every operator with
spatial context learned to predict change and scored far above
persistence. Learning depended on a loss that weights the lesion mask:
trained with a plain mean squared error, the canonical operator learned no
change. A soft-Dice term on the mask added about 0.05 to the change-region
Dice. The interval length has to reach the operator, but the form in which
it does so mattered little: a fourth-order Runge--Kutta step gave no
measurable gain over a single Euler step.

For the second part, the score of an operator did not follow from its
architecture class. The three strongest operators (a finite-element
network, a graph network and a U-Net) come from unrelated classes. At one
year they reach a change-region Dice of about 0.50 to 0.54, against zero
for persistence, and their scores differ by at most about 0.04. The
finite-element network scores highest: above the U-Net at both seeds, and
slightly above the graph network, a lead that is not established. The
finite-element network carries the transport term and the graph network
the extended reach described below. A few properties, each measured one at
a time, separate the operators:
\begin{itemize}
  \item \textbf{Spatial context.} A model without spatial context predicts
    no change; one ring of neighbours raises the change-region Dice from 0
    to about 0.45.
  \item \textbf{Physical reach at the scale of the lesion front.} Placing
    the neighbours of a graph network so that one step reaches as far as
    the front moves between visits raised the change-region Dice by 0.064,
    at an unchanged parameter count.
  \item \textbf{An explicit transport term.} Adding a learned transport
    term to a Finite Element Network raised the change-region Dice by 0.062
    over the same network without the term, and a smaller T-FEN with
    fewer parameters than that network scored the same.
\end{itemize}
Reach and transport are two unrelated ways to correct the same deficit: a
single update that cannot see as far as the disease moves between two
visits. That two independent routes lead to the same level is the central
architectural result of this thesis; it is more informative than the score
of any single model. A third candidate, a local path for the global FNO,
is not established. The FNO scores below the U-Net at one year at both
seeds, and a small local path brought it level at one seed but not at the
second.

Several other changes did not improve the forecast. Larger models did not
help in any of the four classes tested, and neither did a penalty against
lesion shrinkage. The moving-mesh extension, reported in
Appendix~\ref{app:mm-pde}, did not improve accuracy, because its
correction branch never engaged. The validity diagnostics also restrict
how the canonical model may be described: it is a one-step map conditioned
on the step length and does not behave as a continuous-time model. A model
that does behave as one, the energy-conserving T-FEN, can be trained, but
this property does not make it more accurate.

Clinically, the operators predict where the lesion grows better than how
fast an individual eye grows. On where it grows, the growth-region Dice of
the canonical model lies within the range of prior work
(\S\ref{sec:experiments:qualitative}). On how fast, the model reaches a
correlation of 0.40 with the true growth rates, against 0.61 for the model
of \citet{Mai2024}, which starts from the raw OCT volume instead of
segmented masks, 0.52 for \citet{Gigon2021} and about 0.57 to 0.59 for
\citet{Yoshida2025}.

Every conclusion drawn here is limited by the small cohort, the censoring
by the fixed crop and the lack of memory of past visits. Within these
limits, for forecasts from longitudinal medical imaging in a low-data
clinical setting, the framework around the model and the properties of
the operator placed in it mattered more than the choice of operator
family. Of these properties, the most important was whether one update can
see as far as the disease moves between two visits.
% 2026-10-19 rewritten (supervisor review 2, SV-7.1-SV-7.5; AI 6-7.44-6-7.56):
% "requires", "equally", "implements these requirements"; no participle
% opener, no "is not what decides", no "What separates ... are", no slogan
% sentence, no closing "what matters is ... rather than"; the Runge-Kutta
% null once; reach/transport interpretation directly after the list; limits
% in the last paragraph. Numbers unchanged (Table 5.1, 5.3, 5.5, 5.7).



%%
%%%%%%%%%%%%%%%%%%%%%%%%%%%%%%%%%%%%%%%%%%%%%%%%%%%%%%%%%%%%%%%%%%%%%%%%%%%%%%%%
